% page 353 of the facsimile (image p-352 of the scan)
% block 175.3 x 233.3 mm ; left margin 26.7 mm
\pgc{5.080mm}{0.000mm}{
\l{\cel{59}345}
\l{}
\l{\cel{4}\sou{lotuso}. (\sou{bot.}) I. Speco di jujubiero. Che la antiqui : arboro}
\l{\cel{7}di qua la frukto egardesis kom iganta la stranjeri qui kon-}
\l{\cel{7}sumis de lu obliviar sua patrio. II. (\sou{bot.)}\cel{2}Speco di nenu-}
\l{\cel{7}faro blua qua kreskas en Egiptia. - L.\cel{1}\sou{nelumbium speciosum}}
\l{\cel{7}- DEFIRS.}
\l{}
\l{\cel{4}\sou{loxodromo}. I. (\sou{geom.}) Kurvo trasita sur la surfaco di sfero e}
\l{\cel{7}qua intersekas sam-angule omna meridiani. - II. (\sou{maro-navig.})}
\l{\cel{7}Kurvo quan trasas navo kande lu iras kontinue ad un punto di}
\l{\cel{7}la busolo, en un rumbo di vento. - DEFIS.}
\l{}
\l{\cel{4}\sou{loyala}. Skrupuloze fidela ye sua promisi, ye sua su-obligi. -}
\l{\cel{7}DEFIRS.}
\l{}
\l{\cel{4}lu. Pronomo personala, triesma persono, singularo, epicena.}
\l{}
\l{\cel{4}\sou{lubrifikar}. (\sou{trans.)}I.Igar glitigiva (tisuo) por igar plu facila}
\l{\cel{7}lua funciono. - II. Untizar ye graso (mashino, aparato). EFIS.}
\l{}
\l{\cel{4}\sou{lucio}. (\sou{zool.}) Fisho qua vivas en aquo sen-sala, tre devorera,}
\l{\cel{7}ek la familio \textquotedbl{}ezoci\textquotedbl{}, di qua la korpo esas funelatra, kun}
\l{\cel{7}muzelo longa, boko tre fendita e garnisita ye granda quanto}
\l{\cel{7}de denti, e di qua la karno\cel{2}tre prizesas. - L.\cel{1}\sou{esox lucius}.}
\l{\cel{7}- ISL.}
\l{}
\l{\cel{4}\sou{lucida}. Qua komprenas klare la kozi. - EFIS.}
\l{}
\l{\cel{4}\sou{luciolo}. (\sou{zool.)}\cel{2}Lampiro aloza e fosforecanta. - L.\cel{1}\sou{luciola}}
\l{\cel{7}\sou{italica}. - FIS.}
\l{}
\l{\cel{4}\sou{ludar}. (\sou{netrans}.)(\sou{per)}.IAgar ulo por amuzar su. II. Amuzar su}
\l{\cel{7}per manuagar instrumento. - III.Amuzar su, kun altra persono}
\l{\cel{7}o personi, per agi en qua uli ganas e la ceteri perdas. -}
\l{}
\l{}
\l{\cel{4}\sou{ludiono}. (\sou{fiziko}).Figureto, statueto, imajeto, ek vitro, garni-}
\l{\cel{7}sita ye la parto supera ye sfereto vakua qua flotacas segun}
\l{\cel{7}la principo da Arkimedes, ed acensas o decensas (segun ke}
\l{\cel{7}onu lo deziras) en aquo. - FIS.}
\l{}
\l{\cel{4}\sou{lugar}. (\sou{trans.,}\cel{1}ad). (\sou{aludante proprietero}). Cedar la uzo-yuro}
\l{\cel{7}di kozo - cedo dum-frista, po preco konvencionata. -}
\l{}
\l{\cel{4}\sou{lugro}. (\sou{maronavigado)}. Milito-naveto lejera, uzita precipue da}
\l{\cel{7}la kontrabandisti e da la pirati. - DEFIS.}
\l{}
\l{\cel{4}\sou{luko}. (\sou{navo.}) Aperturo quadriangula en la ponto, por igar komuni-}
\l{\cel{7}kar du etaji. - DER.}
\l{}
\l{\cel{4}\sou{lukano.}\cel{2}(\sou{zool}.)\cel{2}Insekto aloza, granda, kun mandibuli di qui}
\l{\cel{7}la formo esas olta di korni dentoza, kelke simila a la kornaco}
\l{\cel{7}di cervo. - L.\cel{1}\sou{lucanus}. - FL.}
}
